\chapter{Boolean Analysis --- Threshold Functions --- Low Degree Norm}
%%% generated by the blueprint pipeline from TCSlib.BooleanAnalysis.ThresholdFunctions.LowDegreeNorm
\section{Overview}
Hypercontractive moment bounds and interpolation compare the first and second moments of low-degree Walsh polynomials with the sharp exponential degree factor.

\section{Declarations}

\begin{lemma}[Low-degree energy retained by noise]
\lean{BooleanAnalysis.ThresholdFunctions.noise_degree_lower}
\leanok
\uses{BooleanAnalysis.Hypercontractivity.noise_l2_fourier, BooleanAnalysis.ThresholdFunctions.MultilinearPolynomial, BooleanAnalysis.ThresholdFunctions.MultilinearPolynomial.HasDegreeAtMost, BooleanAnalysis.ThresholdFunctions.MultilinearPolynomial.eval, BooleanAnalysis.ThresholdFunctions.MultilinearPolynomial.fourierCoeff_eval, BooleanAnalysis.expect, BooleanAnalysis.innerProduct, BooleanAnalysis.noiseOp}
\difficulty{4}
Let $n\ge0$ and give $\{0,1\}^n$ uniform probability; write $\mathbb E$ for expectation. Put $\chi_S(x)=\prod_{i\in S}(-1)^{x_i}$ and $\widehat h(S)=\mathbb E[h\chi_S]$ for $S\subseteq[n]$ and real cube functions $h$. Let $P(x)=\sum_{S\subseteq[n]}a_S\chi_S(x)$ with real coefficients; degree at most $k$ means $a_S=0$ for $|S|>k$. If $P$ has degree at most the integer $k\ge0$ and $0\le\rho\le1$, define $T_\rho P=\sum_S\rho^{|S|}a_S\chi_S$. Then $\rho^{2k}\sum_Sa_S^2\le\mathbb E[(T_\rho P)^2]$, with natural powers including $0^0=1$.
\end{lemma}

\begin{lemma}[Hypercontractive bound for a low-degree moment]
\lean{BooleanAnalysis.ThresholdFunctions.low_degree_moment_bound}
\leanok
\uses{BooleanAnalysis.Hypercontractivity.expect_rpow_abs_nonneg, BooleanAnalysis.Hypercontractivity.expect_sq_noiseOp_nonneg, BooleanAnalysis.Hypercontractivity.general_one_function_hypercontractivity, BooleanAnalysis.ThresholdFunctions.MultilinearPolynomial, BooleanAnalysis.ThresholdFunctions.MultilinearPolynomial.HasDegreeAtMost, BooleanAnalysis.ThresholdFunctions.MultilinearPolynomial.eval, BooleanAnalysis.ThresholdFunctions.noise_degree_lower, BooleanAnalysis.expect, BooleanAnalysis.noiseOp}
\difficulty{5}
Let $n\ge0$ and give $\{0,1\}^n$ uniform probability; write $\mathbb E$ for expectation. Put $\chi_S(x)=\prod_{i\in S}(-1)^{x_i}$ and $\widehat h(S)=\mathbb E[h\chi_S]$ for $S\subseteq[n]$ and real cube functions $h$. Let $P(x)=\sum_{S\subseteq[n]}a_S\chi_S(x)$ with real coefficients; degree at most $k$ means $a_S=0$ for $|S|>k$. If $P$ has degree at most the integer $k\ge0$ and $1<t<2$, then $(t-1)^k\sum_Sa_S^2\le(\mathbb E|P|^t)^{2/t}$.
\end{lemma}

\begin{lemma}[Interpolation between first and second moments]
\lean{BooleanAnalysis.ThresholdFunctions.moment_interpolation}
\leanok
\uses{BooleanAnalysis.BoolCube, BooleanAnalysis.BooleanFunc, BooleanAnalysis.expect, BooleanAnalysis.expect_eq_fintypeExpect}
\difficulty{5}
Let $n\ge0$ and give $\{0,1\}^n$ uniform probability; write $\mathbb E$ for expectation. For every real cube function $f$ and real $1<t<2$, $\mathbb E|f|^t\le(\mathbb E|f|)^{2-t}(\mathbb E[f^2])^{t-1}$.
\end{lemma}

\begin{lemma}[Natural-power form of the interpolation bound]
\lean{BooleanAnalysis.ThresholdFunctions.interp_pow}
\leanok
\difficulty{4}
Let $m\ge0$ be an integer and let $A,B,M\ge0$. If $M\le A^{1/(m+1)}B^{m/(m+1)}$, then $M^{2(m+1)}\le A^2B^{2m}$, with every zero base to the zero power interpreted as one.
\end{lemma}

\begin{lemma}[Natural-power form of a moment comparison]
\lean{BooleanAnalysis.ThresholdFunctions.hc_pow}
\leanok
\difficulty{4}
Let $m,k\ge0$ be integers and let $B,M,D\ge0$. If $D^kB\le M^{2(m+1)/(2m+1)}$, then $(D^kB)^{2m+1}\le M^{2(m+1)}$, using natural powers and $0^0=1$.
\end{lemma}

\begin{lemma}[An algebraic second-moment comparison]
\lean{BooleanAnalysis.ThresholdFunctions.abstract_cancel}
\leanok
\difficulty{5}
Let $m,k\ge0$ be integers and $A,B,M,C,D$ be real numbers with $B>0$, $C\ge1$, and $CD=1$. If $(D^kB)^{2m+1}\le M^{2(m+1)}$ and $M^{2(m+1)}\le A^2B^{2m}$, then $B\le(C^{m+1})^{2k}A^2$.
\end{lemma}

\begin{lemma}[Finite exponential degree factor]
\lean{BooleanAnalysis.ThresholdFunctions.finite_algebra}
\leanok
\uses{BooleanAnalysis.ThresholdFunctions.abstract_cancel, BooleanAnalysis.ThresholdFunctions.hc_pow, BooleanAnalysis.ThresholdFunctions.interp_pow}
\difficulty{4}
Let $m\ge1$ and $k\ge0$ be integers, let $A,M\ge0$ and $B>0$, and put $t=(2m+1)/(m+1)$. If $(t-1)^kB\le M^{2/t}$ and $M\le A^{2-t}B^{t-1}$, then $B\le((1+1/m)^{m+1})^{2k}A^2$.
\end{lemma}

\begin{theorem}[Limit of the finite degree factors]
\lean{BooleanAnalysis.ThresholdFunctions.degree_norm_limit}
\leanok
\difficulty{4}
For every integer $k\ge0$, the sequence $((1+1/m)^{m+1})^{2k}$ converges to $e^{2k}$ as the positive integer $m\to\infty$.
\end{theorem}

\begin{theorem}[Endpoint of a family of norm comparisons]
\lean{BooleanAnalysis.ThresholdFunctions.degree_norm_endpoint}
\leanok
\uses{BooleanAnalysis.ThresholdFunctions.degree_norm_limit}
\difficulty{1}
Let $k\ge0$ be an integer and $A,B\in\mathbb R$. If $B\le((1+1/m)^{m+1})^{2k}A^2$ for every integer $m\ge2$, then $B\le e^{2k}A^2$.
\end{theorem}

\begin{theorem}[Sharp exponential first-to-second moment bound]
\lean{BooleanAnalysis.ThresholdFunctions.low_degree_l1_l2_sq}
\leanok
\uses{BooleanAnalysis.Hypercontractivity.expect_rpow_abs_nonneg, BooleanAnalysis.ThresholdFunctions.MultilinearPolynomial, BooleanAnalysis.ThresholdFunctions.MultilinearPolynomial.HasDegreeAtMost, BooleanAnalysis.ThresholdFunctions.MultilinearPolynomial.eval, BooleanAnalysis.ThresholdFunctions.MultilinearPolynomial.expect_eval_sq, BooleanAnalysis.ThresholdFunctions.degree_norm_endpoint, BooleanAnalysis.ThresholdFunctions.expect_nonneg, BooleanAnalysis.ThresholdFunctions.finite_algebra, BooleanAnalysis.ThresholdFunctions.low_degree_moment_bound, BooleanAnalysis.ThresholdFunctions.moment_interpolation, BooleanAnalysis.expect}
\difficulty{5}
Let $n\ge0$ and give $\{0,1\}^n$ uniform probability; write $\mathbb E$ for expectation. Put $\chi_S(x)=\prod_{i\in S}(-1)^{x_i}$ and $\widehat h(S)=\mathbb E[h\chi_S]$ for $S\subseteq[n]$ and real cube functions $h$. Let $P(x)=\sum_{S\subseteq[n]}a_S\chi_S(x)$ with real coefficients; degree at most $k$ means $a_S=0$ for $|S|>k$. If $P$ has degree at most the integer $k\ge0$, then $e^{-2k}\sum_{S\subseteq[n]}a_S^2\le(\mathbb E|P|)^2$.
\end{theorem}
